\documentclass[10pt]{article}
\usepackage{lmodern}
\usepackage[T1]{fontenc}
\usepackage[utf8]{inputenc}
\usepackage{microtype}
\usepackage{enumitem}
\usepackage{bm}
\usepackage{amsmath}
\usepackage{amssymb}
\usepackage[bookmarks=false,breaklinks=false,pdfborder={0 0 1},backref=false,colorlinks=false]{hyperref}
\usepackage{graphicx}
\usepackage{xcolor}
\usepackage{array}
\usepackage{booktabs}

\oddsidemargin -0.04cm
\evensidemargin -0.04cm
\textwidth 16.59cm
\textheight 21.94cm
\topmargin -0.25cm
\renewcommand{\baselinestretch}{1.35}
\setlength{\parindent}{0pt}
\setlength{\parskip}{0.55em plus 0.1em minus 0.1em}

\newcommand{\paperTitle}{SCOPE-FD: A Client Selection Method in Federated Distillation for Massive-MIMO Systems}

\begin{document}

\title{
\small Response to Reviewers' Comments and Changes Incorporated in the Revised Manuscript \\
\vspace{1em}
{\Large \bf \paperTitle} \\
(\textbf{Original Manuscript ID: TAI-2026-May-A-00956})
}
\date{}
\maketitle

\vspace{-3em}

We sincerely thank the Editor, the Associate Editor, and all the anonymous Reviewers for their valuable time and effort in evaluating our manuscript. We are also thankful to all of you for the detailed and constructive feedback. The comments were extremely valuable and have helped us improve the quality, the clarity, and the technical depth of our work. In the revised version, we have carefully addressed the raised concerns. {\color{blue}All changes made in response to the Reviewers' comments have been clearly highlighted in blue, as per the journal's guidelines}. Below, we provide a detailed point-by-point response to each comment. Please note that the Reviewers' comments are \textit{italicized}, followed by our responses in normal font. Along with this letter we submit (a) a clean version of the revised manuscript and (b) a coloured marked version of the revised manuscript, in which the changes are highlighted.\\

\hspace{0.75cm}{\bf {\Large {~~~~~~~~~~~~Response to Comments of Associate Editor}}}

\noindent
\textbf{Comment:} \textit{A revision stage is needed.}\\
\noindent
\textbf{Reply:} {We thank the Associate Editor for the opportunity to revise and resubmit the manuscript. The revision covers the following. Every experimental claim is now supported by multiple seeds with standard deviations. We added a four-way ablation separating the participation-debt term from the under-prediction and coverage terms, and four published client-selection baselines ported into the federated distillation pipeline. We added a full grid over the two scoring coefficients, extended the heterogeneity and client-scale sweeps, included configurations in which $K$ does not divide $N$, and added a rolling-window fairness metric that does not align with the participation cycle. Section~V-B has been replaced with a convergence result for the partial-participation setting. The revised manuscript runs to thirteen pages. The figures have been consolidated into three multi-panel floats, so that the sweep, ablation and robustness studies each read as one figure, and no result has been dropped in the process. The changes are highlighted in blue in the revised manuscript, and each point is addressed in the per-Reviewer responses below.}\\

\hspace{0.75cm}{\bf {\Large {~~~~~~~~~~~~Response to Comments of Reviewer 1}}}

\noindent
\textbf{Comment 1:} \textit{The paper aims to demonstrate convergence compatibility with FedTSKD and client selection. However, the convergence proof is based on FedTSKD, and the effect of client selection via a client subset with bounded variance can easily be derived. Therefore, the analytical contribution of Section V-B is minimal.} \\
\noindent
\textbf{Reply 1:} {We thank the Reviewer for this comment. We have replaced Section~V-B with a convergence result for the partial-participation setting.

The former Section~V-B observed that the per-client bound of [13] continues to hold under our selector, which, as the Reviewer notes, follows immediately: [13] has no client-selection mechanism, so no step of its analysis was conditioned on one. The substantive question is different. The index $t$ in that bound counts the local-training and distillation steps a client has actually taken. Under full participation it advances with the round counter, but under partial participation a client's clock advances only in the rounds where it is selected, so the bound of [13] alone says nothing in terms of the communication budget $R$. Closing that gap requires a lower bound on how often each client is served, which is what our Proposition~1 provides.

Section~V-B now composes the two. Corollary~1 gives a per-client rate of $\mathcal{O}(N/(KR))$ in communication rounds that holds for every client at every finite horizon rather than on average, and Corollary~2 carries the same rate to the global objective. Both recover the $\mathcal{O}(1/R)$ rate of [13] when the cohort is the whole pool. No assumption beyond those of [13] is required. Please refer to Section~V-B of the revised manuscript.}\\

\noindent
\textbf{Comment 2:} \textit{Only two datasets, Fashion-MNIST and MNIST, which have very similar characteristics, are tested. The dataset is compatible with the given model. However, as mentioned in the conclusion and future work, the applicability of the proposed method to domains other than images is questionable, as the proposed SCOPE-FD was tested under a narrow scenario.} \\
\noindent
\textbf{Reply 2:} {We thank the Reviewer for this comment. We have evaluated the selector on a domain outside images and completed the MNIST study.

On the Free Spoken Digit Dataset, a ten-class corpus of spoken digits represented as log-spectrograms and paired with an \texttt{AudioCNN} model, the proposed selector reaches $46.16 \pm 1.69$ over five seeds, against $45.54 \pm 1.17$ for the debt-only variant, $44.99 \pm 1.59$ for uniform random, $40.66 \pm 1.48$ for SubTrunc and $40.53 \pm 1.49$ for DivFL. The participation Gini is $1.33\%$ against $12.49\%$ for uniform random and about $45\%$ for the two submodular methods. The proposed selector leads uniform random on every one of the five seeds. Two settings were adjusted for the size of this corpus, namely three local epochs in place of one and a public set of $150$ samples, and both are stated in the manuscript.

On MNIST over three seeds, the debt-only variant reaches $90.88 \pm 0.42$, uniform random $90.58 \pm 0.43$ and the complete score $90.50 \pm 0.56$, with DivFL at $88.48 \pm 0.43$ and SubTrunc at $87.91 \pm 0.50$.

The reading is consistent across the three datasets. The fairness separation is identical everywhere, since it is fixed by the participation arithmetic, and the accuracy ordering is preserved. What changes is the absolute accuracy level, which reflects task difficulty. We also removed the earlier single-seed EMNIST replication, which these multi-seed studies supersede. Please refer to Section~VI-L of the revised manuscript.}\\

\noindent
\textbf{Comment 3:} \textit{As the experimental results were conducted in limited settings, it would be useful to show confidence intervals on the graphs.} \\
\noindent
\textbf{Reply 3:} {We thank the Reviewer for this suggestion. Every result is now reported over multiple independent seeds, five for the main configurations and three for the sweeps, with the mean and the standard deviation for every entry and the number of seeds stated in each table. Every learning curve carries a shaded band and every point estimate a whisker, both at one standard deviation over seeds. As an example, the headline configuration at $N=30$, $K=5$, $\alpha=0.5$ now reads $71.21 \pm 0.99$ for SCOPE-FD and $70.99 \pm 1.38$ for uniform random. The reporting protocol is in Section~VI-A, and the multi-seed values appear in Table~II and in Fig.~2 and Fig.~3 of the revised manuscript.}\\

\noindent
\textbf{Comment 4:} \textit{The coefficients (e.g., alphau and alphad) are not ablated. The selection process for these coefficients needs to be described.} \\
\noindent
\textbf{Reply 4:} {We thank the Reviewer for this suggestion. We have added a full grid over $\alpha_u \in \{0, 0.1, 0.2, 0.3, 0.5, 0.8\}$ and $\alpha_d \in \{0, 0.05, 0.1, 0.2, 0.4\}$, giving $30$ cells with three seeds each. Final accuracy across the whole grid ranges from $70.94\%$ to $72.27\%$, a spread of $1.33$ percentage points, and the participation Gini is $1.33\%$ in every cell. This matches the magnitude argument of Section~IV-D: the debt term dominates the score whenever $\alpha_u + \alpha_d < 1$, so the fairness guarantee holds for any admissible choice. The setting used throughout the paper, $(\alpha_u, \alpha_d) = (0.3, 0.1)$, reaches $71.99\%$, within $0.27$ percentage points of the best cell. The selection of these values is now described in a dedicated paragraph, ``Choice of the coefficients'', in Section~IV-D rather than in a footnote. The grid is reported in full in Section~VI-C of the revised manuscript.}\\

\noindent
\textbf{Comment 5:} \textit{The baseline is too simple when considering uniform randomization. There are other client selection mechanisms in FL (e.g., DevFL) that apply submodular diversity maximization to the client histogram, which is similar to SCOPE-FD's coverage penalty. Therefore, the authors need to consider more realistic and fair comparison baselines.} \\
\noindent
\textbf{Reply 5:} {We thank the Reviewer for this suggestion. We have added four published client-selection methods, each ported into the federated distillation pipeline. DivFL performs facility-location submodular maximization; since a client in federated distillation does not send a gradient, we use the public-set logit vector as the representation for the submodular objective. SubTrunc is a truncated submodular method diversified by client loss. UnionFL is diversified by historical participation. Oort is a utility-based selector from standard federated learning, run without modification.

At the headline configuration over five seeds, SCOPE-FD reaches $71.21 \pm 0.99$ with a Gini of $1.33\%$, uniform random $70.99 \pm 1.38$ with $12.49\%$, UnionFL $71.18 \pm 1.40$ with $1.33\%$, DivFL $69.88 \pm 0.83$ with $36.84\%$, SubTrunc $69.40 \pm 1.00$ with $38.65\%$ and Oort $21.14 \pm 0.60$ with $83.33\%$. SCOPE-FD leads DivFL by $1.33$ and SubTrunc by $1.81$ percentage points in accuracy and lowers their participation Gini by more than an order of magnitude. Against uniform random and UnionFL the accuracy difference lies inside one standard deviation, so the three are tied on that axis. UnionFL also returns $1.33\%$ and the same $13.0$ rounds here, so it matches the proposed selector on every column of Table~II and we now separate the two elsewhere rather than claiming a fairness margin that the table does not show. Over the cohort sweep at $N=30$ the proposed selector returns $6.67\%$, $0.00\%$, $1.33\%$ and $0.67\%$ at $K \in \{1,3,5,10\}$, each the value the closed form specifies, at a seed standard deviation of zero. UnionFL returns $55.56\%$, $29.44\%$, $1.33\%$ and $22.89\%$ on the same runs, its seed standard deviation reaching $3.11$. Its agreement at $K=5$ is therefore a coincidence of one configuration. The ported methods are described in Section~II-A and the results are in Section~VI-B and Table~II of the revised manuscript.}\\

\noindent
\textbf{Comment 6:} \textit{Only a single value of Dirichlet $\alpha$ (0.5) is tested, producing moderate heterogeneity. Experiment scale is too small to represent Massive-MIMO systems.} \\
\noindent
\textbf{Reply 6:} {We thank the Reviewer for raising both points, and we have addressed each with new experiments.

For heterogeneity, the Dirichlet parameter is now swept over $\alpha \in \{0.05, 0.1, 0.3, 0.5, 1.0, 5.0\}$ together with an independent and identically distributed setting. At $\alpha = 0.05$ the proposed selector reaches $26.74 \pm 4.19$ against $27.07 \pm 4.24$ for uniform random, and at $\alpha = 5.0$ it reaches $84.67 \pm 0.11$ against $84.60 \pm 0.18$, so accuracy is tied at every level. The participation Gini is $1.33\%$ for the proposed selector at every level, against $12.16\%$ to $12.64\%$ for uniform random, so the fairness advantage is invariant to the degree of heterogeneity.

For scale, the revised paper reports client pools of $N \in \{30, 47, 50, 53, 100, 200\}$ with cohorts of $K$ up to $40$, spanning participation ratios from $3.3$ to $33$ percent, all over five seeds. The fairness law holds at every pool size. At $N = 200$ with $K = 10$ the proposed selector reaches $51.27\%$ accuracy with a cycle-aligned Gini of $0.00\%$ against $48.84\%$ and $24.16\%$ for uniform random, and at $N = 200$ with $K = 20$ and $K = 40$ the Gini remains $0.00\%$ against $16.24\%$ and $10.53\%$. Those zeros are an alignment of the horizon rather than an effect of the pool, since $KR$ is an exact multiple of $N$ at $R = 100$, and the revised Section~VI-K now says so rather than reading them as a pool-size trend.

We would add that the quantity which sets the scale of a scheduler on an mMIMO backhaul is the per-round cohort rather than the size of the pool behind it. Recovering the $K$ per-client logit streams by zero forcing requires the composite uplink channel to have full column rank, so $K \le N_{BS}$, and at the $N_{BS} = 64$ antennas of our setup the largest cohort we report, $K = 40$, already occupies $40$ of the $64$ spatial dimensions the array can separate. The sweep therefore covers most of the range the physical layer admits, and this bound is now stated explicitly in Section~III-A and related to the sweep in Section~VI-K.

The enlarged study also lets us characterize where the accuracy advantage holds. Across fifteen matched configurations spanning $N$ from $30$ to $200$ and participation ratios from $3.3$ to $33$ percent, the accuracy gap over uniform random falls as the ratio grows, at a Spearman rank correlation of $-0.815$, from $+2.43$ percentage points at $N = 200$, $K = 10$ to $-1.42$ at $N = 200$, $K = 40$. Twelve of the fifteen gaps are positive, and the three that are not all use a cohort of at least twenty clients, where the debt-only variant leads the complete score. The fifteen configurations share a seed set, so the correlation is now reported as a trend and carries no significance level. The manuscript now confines the accuracy claim to the sparse regime and notes that a deployment at a high participation ratio should prefer the debt-only variant. The heterogeneity sweep is in Section~VI-D and Fig.~2(a) and (b), the scaling study in Section~VI-K and Fig.~3(b), and the ratio analysis in Fig.~3(c) of the revised manuscript.}\\

\hspace{0.75cm}{\bf {\Large {~~~~~~~~~~~~Response to Comments of Reviewer 2}}}

\noindent
\textbf{Comment 1:} \textit{The most important missing comparison is a debt-only selector with $\alpha_u = \alpha_d = 0$. Section IV-A states that this setting reduces SCOPE-FD to deterministic round-robin, and Section V-A shows that the debt term dominates the fairness guarantee. Without this baseline, it is unclear whether the gains in Fig. 3 come from the under-prediction and coverage terms or simply from replacing random selection with balanced rotation. Please also report ablations for debt plus under-prediction only and debt plus coverage only.} \\
\noindent
\textbf{Reply 1:} {We thank the Reviewer for this comment. We have added the debt-only selector and the two partial variants, and the complete four-way ablation now runs at the headline configuration, across the participation sweep and across the channel sweep, with five seeds.

At the headline configuration the debt-only selector reaches $71.26 \pm 1.14$, debt plus under-prediction $70.55 \pm 1.19$, debt plus coverage $71.56 \pm 0.98$ and the complete score $71.21 \pm 0.99$, with a participation Gini of $1.33\%$ for all four. The four variants are statistically indistinguishable in final accuracy, and the fairness guarantee comes entirely from the participation-debt term. The measurable contribution of the two auxiliary terms is to convergence speed: the complete score reaches sixty percent absolute accuracy in $13.0$ rounds against $14.0$ for debt only, and uniform random needs $19.0$.

We have accordingly rewritten the claims in the Abstract, in Section~I-A and in Section~VII so that the deterministic fairness guarantee is presented as the main result and the two auxiliary terms as a convergence refinement. The ablation is in Section~VI-B and Table~II of the revised manuscript.}\\

\noindent
\textbf{Comment 2:} \textit{Equations (12) and (14) require each client to reveal its label histogram $h_i$. This is not raw-data sharing, but it still discloses the local class distribution. Since FD is partly motivated by privacy, this assumption should be treated as a central limitation. Please either evaluate the Laplace-noise or server-side surrogate variants mentioned in Section IV-D, or narrow the privacy-related claims.} \\
\noindent
\textbf{Reply 2:} {We thank the Reviewer for this comment. We have evaluated both variants named in Section~IV-D. The first applies a Laplace mechanism to the histogram before the selector reads it, swept over a privacy budget $\varepsilon_{\mathrm{dp}} \in \{0.1, 0.5, 1, 2, 5\}$. The second removes the disclosure altogether by replacing the histogram with a server-side surrogate built from the public-set logits each client already submits.

Neither carries a measurable cost. Against $71.99 \pm 0.63$ for the undefended selector, the Laplace variants reach $72.01 \pm 0.86$ at $\varepsilon_{\mathrm{dp}} = 0.1$ and $71.85 \pm 0.73$ at $\varepsilon_{\mathrm{dp}} = 5$, and the surrogate reaches $71.93 \pm 0.67$. Every difference lies inside one standard deviation, including at the strongest noise setting, and the participation Gini is $1.33\%$ for all seven variants.

The reason is structural. The histogram feeds only the under-prediction and coverage terms, which the ablation of Comment~1 shows do not change final accuracy, while the participation-debt term that carries the fairness guarantee never reads the histogram. Since the disclosure can be removed at no measured cost, the privacy-related claims did not need narrowing. Please refer to Section~VI-H of the revised manuscript.}\\

\noindent
\textbf{Comment 3:} \textit{The participation-fairness proof is convincing, but the convergence claim needs more care. Section V-B says that the FedTSKD $O(1/t)$ bound holds ``without modification,'' but SCOPE-FD changes full participation into partial participation. Please provide a clearer theorem for the partial-participation setting, or state more modestly that SCOPE-FD is compatible with the FedTSKD pipeline under the assumptions of [13].} \\
\noindent
\textbf{Reply 3:} {We thank the Reviewer for this comment. We have taken the first of the two options offered and provided the theorem.

We first verified which assumptions of [13] are at stake. Its Assumptions~1--4 are smoothness and strong convexity of the local training and distillation losses together with bounded variance and bounded norm of their stochastic gradients, and none of the four refers to the selected subset. The phrase ``uniformly and independently'' in Assumption~3 describes the minibatch drawn inside a client rather than the way clients are drawn. Its Algorithm~1 admits all clients in every round and its Section~VI adopts full participation by design, so there is no client-sampling assumption in [13] for a deterministic rotation to reconcile with.

What partial participation does change is the time index. In [13] the index $t$ counts the steps a client has actually taken, which coincides with the round counter only under full participation; under our selector a client's clock advances only in the rounds where it is selected. Section~V-B now closes that gap. Corollary~1 composes the per-client bound of [13] with the participation-balance guarantee of Proposition~1 and gives a rate of $\mathcal{O}(N/(KR))$ in communication rounds that holds for every client at every finite horizon rather than on average, and Corollary~2 carries it to the global objective. Setting $K=N$ returns the $\mathcal{O}(1/R)$ rate of [13]. Section~V-B also works through the four assumptions one at a time and states which part of the pipeline each constrains. Please refer to Section~V-B of the revised manuscript.}\\

\noindent
\textbf{Comment 4:} \textit{The experiments use a single reported seed and do not show confidence intervals or standard deviations. It remains unclear what the actual performance is under all possible scenarios. Claims such as ``statistically tied'' and ``3x faster'' should be supported by multi-seed results and error bars. The rounds-to-80\%-of-final-accuracy metric is useful, but it should be complemented with final accuracy tables, rounds to a fixed absolute accuracy, and sensitivity to $\alpha_u$, $\alpha_d$, Dirichlet $\alpha$, and K/N.} \\
\noindent
\textbf{Reply 4:} {We thank the Reviewer for this detailed comment, and we have acted on every part of it.

All results now use multiple seeds, five for the main configurations and three for the sweeps. Every table entry is a mean and a standard deviation, every figure carries confidence bands or error bars, and the number of seeds is stated in each table.

We added both metrics requested. The paper reports a plain final-accuracy table, and it replaces the relative rounds-to-eighty-percent-of-final metric with rounds to a fixed absolute accuracy. The threshold reported in Table~II is sixty percent, which every seed of every selector except Oort attains, so the column is a mean over the full set rather than over the subset that happened to cross it. At the stricter seventy percent threshold not every seed finishes within the hundred-round budget, so we report attainment counts there instead of a mean.

We added the four sensitivity studies: a grid of $\alpha_u$ and $\alpha_d$ with $30$ cells, Dirichlet $\alpha$ from $0.05$ to $5.0$ together with the independent and identically distributed case, and participation covering $K \in \{1, 3, 5, 10\}$ at $N = 30$ and ratios of five, ten and twenty percent at $N$ up to $200$.

The two claims the Reviewer names have been revisited. The tie claim is now supported by paired Wilcoxon signed-rank tests over matched seeds, and the speed claim has been replaced with the measured rounds-to-threshold values and their standard deviations. The protocol is in Section~VI-A and the sensitivity studies are in Sections~VI-C, VI-D, VI-E and VI-G of the revised manuscript.}\\

\noindent
\textbf{Comment 5:} \textit{The motivation emphasizes massive-MIMO, energy budgets, and contention-limited uplinks, but the selection score in Eq. (8) does not include any channel or energy information. The SNR sweep mainly shows that FedTSKD remains robust under noise; it does not show that SCOPE-FD exploits MIMO properties. Please either add channel/energy-aware experiments or clearly position SCOPE-FD as a fairness/data-coverage selector running on top of an mMIMO communication substrate.} \\
\noindent
\textbf{Reply 5:} {We thank the Reviewer for this comment. We built the channel-aware and energy-aware variant described, adding a fourth score term that reads per-client channel quality and energy, and evaluated it under an explicit per-round energy budget.

The variant is worse on both axes at every downlink SNR tested. It reaches $61.56 \pm 3.31$ at $-30$~dB against $65.23 \pm 1.39$ for the plain selector, a deficit of about $3.5$ percentage points that is stable across the sweep, and its participation Gini rises to $11.69\%$ against $1.33\%$. Instrumenting the selection loop identifies the mechanism: the energy budget acts as a hard feasibility filter that removes clients whose energy exceeds the remaining allowance regardless of how far behind they are in participation. Under the plain selector every round admits the full cohort of five; under the channel-aware variant $89$ of $100$ rounds admit fewer, the mean cohort falls to $3.56$, the counts span $15$, and two of the thirty clients are never selected. This marks a boundary of Proposition~1, which assumes every client remains selectable.

We have therefore taken the second option the Reviewer offers and now position SCOPE-FD explicitly as a fairness and data-coverage selector operating on top of an mMIMO substrate rather than one that exploits channel state, with the positioning supported by this measurement. Please refer to Section~VI-J of the revised manuscript and to the revised framing in the Abstract and Section~VII.}\\

\noindent
\textbf{Comment 6:} \textit{Please define ``K$|$N'' clearly as a divisibility condition in the abstract. In Section VI-A, fix missing spaces such as ``[13].FMNIST'' and ``sweeps. In every configuration.'' The title should likely use lowercase ``for'' if title case is applied consistently.} \\
\noindent
\textbf{Reply 6:} {We thank the Reviewer for catching these, and all three have been corrected. Rather than define the notation in the abstract, we removed it: the abstract now states the fairness result in words, namely that the participation imbalance depends only on the pool size, the number served per round and the training horizon. The divisibility condition is stated precisely in Section~V-A, where we also note that the exact condition is that $N$ divides $KR$ rather than that $K$ divides $N$. The missing spaces in Section~VI-A have been fixed, including ``[13]. FMNIST'' and ``sweeps. In every configuration.'' The title now uses lowercase ``for''. We also made a further proofreading pass over the whole manuscript for similar spacing errors.}\\

\hspace{0.75cm}{\bf {\Large {~~~~~~~~~~~~Response to Comments of Reviewer 3}}}

\noindent
\textbf{Comment 1:} \textit{The authors clearly explain why traditional federated learning client selection methods cannot be directly applied to federated distillation.} \\
\noindent
\textbf{Reply 1:} {We thank the Reviewer for this observation. We have strengthened that argument with direct evidence rather than leaving it in prose. We ran Oort, a standard utility-based selector from federated learning, unmodified inside the federated distillation pipeline at the headline configuration. It reaches $21.14 \pm 0.60$ accuracy with a participation Gini of $83.33\%$, against $71.21 \pm 0.99$ and $1.33\%$ for SCOPE-FD and $70.99 \pm 1.38$ for uniform random.

The mechanism is the aggregation rule. In federated distillation the server averages logits on the public set with weights proportional to client data size. A utility-based selector repeatedly picks the same small group of high-loss clients, which in federated learning is useful because those updates enter the parameter average directly, but in federated distillation narrows the set contributing to the consensus logits so that the distillation target collapses toward a small subset. The port is described in Section~II-A and the measurement is in Section~VI-B and Table~II of the revised manuscript.}\\

\noindent
\textbf{Comment 2:} \textit{How sensitive is the proposed algorithm to different values of $\alpha_u$ and $\alpha_d$?} \\
\noindent
\textbf{Reply 2:} {We thank the Reviewer for this question. The algorithm is not sensitive to these two values. We ran a grid over $\alpha_u \in \{0, 0.1, 0.2, 0.3, 0.5, 0.8\}$ and $\alpha_d \in \{0, 0.05, 0.1, 0.2, 0.4\}$, giving $30$ cells with three seeds each. Final accuracy across the whole grid stays between $70.94\%$ and $72.27\%$, a spread of $1.33$ percentage points, and the participation Gini is $1.33\%$ in every cell. This is the behavior predicted in Section~IV-D, since the debt term dominates the score whenever $\alpha_u + \alpha_d < 1$ and the fairness guarantee therefore holds for any admissible choice. Please refer to Section~VI-C of the revised manuscript.}\\

\noindent
\textbf{Comment 3:} \textit{Improve the resolution of Figure 1.} \\
\noindent
\textbf{Reply 3:} {We thank the Reviewer for pointing this out. Figure~1 has been regenerated in vector form so that it stays sharp at any zoom level, and we took the same opportunity to enlarge the fonts, thicken the plot lines and clarify the axis labels and the legend. The same treatment has been applied to every other figure in the manuscript for consistency.}\\

\noindent
\textbf{Comment 4:} \textit{How does SCOPE-FD perform under more severe non-IID settings?} \\
\noindent
\textbf{Reply 4:} {We thank the Reviewer for this question. The heterogeneity sweep now covers Dirichlet $\alpha \in \{0.05, 0.1, 0.3, 0.5, 1.0, 5.0\}$ and the independent and identically distributed case. Severe heterogeneity lowers accuracy for every method alike: at $\alpha = 0.05$ SCOPE-FD reaches $26.74 \pm 4.19$ against $27.07 \pm 4.24$ for uniform random, and at $\alpha = 0.1$ it reaches $39.28 \pm 2.55$ against $39.14 \pm 2.69$, so the two are statistically tied on accuracy at every severity level tested. The difference is in fairness and it is stable, with a participation Gini of $1.33\%$ for SCOPE-FD at every value of $\alpha$ against $12.16\%$ to $12.64\%$ for uniform random. The fairness guarantee therefore does not degrade as the data becomes more heterogeneous. Please refer to Section~VI-D and Fig.~2(a) and (b) of the revised manuscript.}\\

\noindent
\textbf{Comment 5:} \textit{How would the method perform under asynchronous federated learning or client dropout scenarios?} \\
\noindent
\textbf{Reply 5:} {We thank the Reviewer for this question, and we have added both settings. Under dropout each selected client fails to return its logits with probability $p_{\text{drop}}$ and the server aggregates over whoever responded. Under bounded staleness a client may return its logits up to $W$ rounds late, aggregated with a staleness discount.

Under dropout the method retains both advantages. At $p_{\text{drop}} = 0.1$ it reaches $71.16 \pm 0.69$ against $70.95 \pm 1.02$ for uniform random, and at $p_{\text{drop}} = 0.3$ it reaches $69.20 \pm 1.03$ against $69.16 \pm 1.48$. Accuracy falls for both as dropout rises, since fewer clients contribute to each round. The participation Gini moves from $0.80\%$ to $1.88\%$ while uniform random moves from $13.42\%$ to $15.45\%$; dropout is the only perturbation in the revision that moves the Gini at all, because it is the only one that changes which clients actually participated.

Under bounded staleness the guarantee is untouched. The Gini is $1.33\%$ at both $W = 1$ and $W = 2$, identical to the synchronous value, because a late arrival changes when a contribution is aggregated but not whether the client was selected. Accuracy is $71.71 \pm 0.61$ and $71.54 \pm 0.60$ respectively. Please refer to Section~VI-I of the revised manuscript.}\\

\noindent
\textbf{Comment 6:} \textit{Since the proposed method relies on a public dataset, how sensitive is its performance to the quality or distribution of that public dataset?} \\
\noindent
\textbf{Reply 6:} {We thank the Reviewer for this question. We varied the public set along three axes, namely its identity, its size and the level of injected label noise.

Identity matters most: replacing the held-out public set with MNIST lowers accuracy from $71.75 \pm 0.81$ to $67.10 \pm 1.05$. Size matters up to a threshold, with a reduction from $2000$ to $500$ samples costing about $2.5$ percentage points and a further reduction to $100$ costing only $0.17$ more, so a few hundred public samples suffice in this setting. Label noise on the public set has no effect at all, which is structural rather than empirical: the aggregation rule consumes public-set logits and never reads public-set labels, and the measured values at noise levels $0.1$ and $0.3$ match the noise-free baseline exactly.

While preparing this study we found that the manuscript described the default public set incorrectly, and we have corrected it. The text stated that MNIST was used as the public set for the Fashion-MNIST experiments, whereas the experiments use $2000$ unlabeled samples drawn once from the held-out split of the private dataset, which is the convention of the underlying framework; MNIST is a variant in this sensitivity study rather than the default. Table~I and Section~VI-A now state this correctly. Since the MNIST cell draws its public set from a separate corpus, its inputs are disjoint from the evaluation data, and the ordering of the selectors is unchanged there, with the proposed selector at $67.10 \pm 1.05$, the debt-only variant at $67.14 \pm 0.95$ and uniform random at $66.72 \pm 0.89$. The absolute accuracy level depends on the choice of public set, but the comparison between selectors does not. Please refer to Section~VI-G of the revised manuscript.}\\

\hspace{0.75cm}{\bf {\Large {~~~~~~~~~~~~Response to Comments of Reviewer 4}}}

\noindent
\textbf{Comment:} \textit{(There are no comments.)} \\
\noindent
\textbf{Reply:} {We thank Reviewer 4 for taking the time to evaluate our manuscript.}\\

\hspace{0.75cm}{\bf {\Large {~~~~~~~~~~~~Response to Comments of the Attached Review Document}}}

\noindent
\textbf{Comment 1 (Convergence analysis):} \textit{The manuscript states that the convergence behavior of FedTSKD is preserved because SCOPE-FD only changes the composition of the selected client subset while leaving the remaining learning procedure unchanged. Since deterministic partial participation introduces a different client-selection policy, it would be helpful to clarify whether the assumptions required in the original convergence analysis of FedTSKD remain satisfied. If a complete theoretical extension is beyond the scope of this work, a brief discussion of the underlying assumptions would improve the presentation.} \\
\noindent
\textbf{Reply 1:} {We thank the Reviewer for this comment. Section~V-B now works through the assumptions of [13] one at a time.

Assumptions~1 and~2 are smoothness and strong convexity of the local training loss and the distillation loss. They are properties of those loss functions alone, and our selector changes no loss function, so they are unaffected. Assumptions~3 and~4 are bounded variance and bounded gradient norm under within-client minibatch sampling; the phrase ``uniformly and independently'' in Assumption~3 describes the minibatch drawn inside a client's own dataset rather than how clients are drawn, and our selector never touches minibatch sampling. The aggregation step reads the selected subset only through the data-size-weighted logit averaging rule, which we use exactly as [13] defines it.

Most directly, there is no client-sampling assumption in [13] for a deterministic rotation either to satisfy or to violate: its Algorithm~1 admits all $N$ clients in every round, its Section~VI states that full participation is adopted by design, and its proof of Theorem~1 tracks one client through its own weight sequence without reading the state of any other client.

Since the audit showed a full extension to be within reach, we have provided one rather than the brief discussion the Reviewer offers as sufficient. What partial participation changes is not an assumption but the time index $t$, which counts the steps a client has actually taken and coincides with the round counter only under full participation. Our Proposition~1 lower-bounds how often each client is served, and composing the two gives Corollary~1 and Corollary~2 in Section~V-B, with a rate of $\mathcal{O}(N/(KR))$ in communication rounds that holds for every client and returns the rate of [13] at $K=N$. Please refer to Section~V-B of the revised manuscript.}\\

\noindent
\textbf{Comment 2 (Fairness evaluation):} \textit{The theoretical participation-fairness guarantee is one of the main strengths of the manuscript. However, the experimental validation could be expanded to demonstrate the proposed behavior under a broader range of settings. For example, additional experiments using different values of N and R, configurations where K does not divide N, or evaluation windows that do not exactly coincide with complete participation cycles would provide stronger empirical evidence for the generality of the proposed fairness behavior.} \\
\noindent
\textbf{Reply 2:} {We thank the Reviewer for this suggestion, and we have added all three stress tests.

For different client-pool sizes, the revised paper covers $N \in \{30, 47, 50, 53, 100, 200\}$ at several participation ratios, and the fairness metrics are now reported at $R \in \{25, 50, 75, 100\}$ rather than only at the final round.

For the case where $K$ does not divide $N$, we added dedicated configurations. At $N = 47$ with $K = 6$ the proposed selector reaches a Gini of $1.40\%$ against $15.70\%$ for uniform random, and at $N = 53$ with $K = 7$ it reaches $1.25\%$ against $13.81\%$. The small residual arises because the total number of selections is not an exact multiple of the pool size, so the counts cannot all be equal.

For evaluation windows that do not align with a participation cycle, we added a rolling-window Gini computed over a window deliberately offset from the cycle length. The proposed selector stays between $4.52\%$ and $18.58\%$ across all tested configurations while uniform random stays between $43.85\%$ and $52.22\%$, so the advantage does not depend on measuring at a cycle boundary.

Working through these tests led to a stronger result that answers all three parts at once. The proof of Proposition~1 bounds the Gini coefficient by replacing an exact numerator with its maximum; retaining that numerator gives the cumulative participation Gini in closed form,
\[
G^{(R)} = \frac{m\,(N-m)}{N K R}, \qquad m = KR \bmod N ,
\]
which requires no new assumption and follows from the same induction. It gives the value at any $N$, at any $R$, and whether or not $K$ divides $N$. It also refines the divisibility statement: the coefficient vanishes exactly when $N$ divides $KR$, not wherever $K$ divides $N$, and our data contains counterexamples in both directions. At the headline setting $N=30$, $K=5$, $R=100$ the cohort divides the pool yet the coefficient is $1.33\%$, while at $N=50$ with $K=3$ the cohort does not divide the pool yet the coefficient is exactly zero. The statement and the caption of Fig.~3(b) have been corrected accordingly. The closed form also makes the horizon dependence explicit rather than monotone, since the coefficient falls to zero whenever the horizon completes a whole number of rotations and rises again between those points, which is why the revised Section~VI-K no longer reads the zeros at the largest pool as a pool-size effect.

Following the Reviewer's point that the fairness behaviour should be shown more broadly, we also report what balanced participation does for the clients rather than only for the counts. The standard deviation of final accuracy across the $N$ clients falls at $K = 1$ from $13.94 \pm 1.24$ under uniform random to $9.11 \pm 2.00$ under the proposed selector, in that direction on all three seeds, and the two converge as the cohort grows, at $6.03$ against $5.99$ by $K = 10$. This is the empirical counterpart of the per-client statement in Corollary~1 and it is now in Section~VI-E of the revised manuscript.

We then tested the expression against the whole campaign. It predicts the measured value correctly in $459$ of $459$ selector runs with no deviation at all, across every dataset, heterogeneity level, SNR, privacy mechanism, staleness window and public-set variation. The only runs it does not cover are the eighteen dropout runs, where a selected client may fail to return so that realized participation differs from intended participation. The closed form and its consequences are in Section~V-A, and the campaign-wide test is in Section~VI-M of the revised manuscript.}\\

\noindent
\textbf{Comment 3 (Experimental reproducibility):} \textit{The reported results appear to be obtained using a single random seed. Since federated learning performance can vary due to random initialization and data partitioning, reporting results over multiple independent runs together with statistics such as mean standard deviation would substantially improve the reproducibility and reliability of the reported performance, particularly for convergence-speed comparisons and the sparse participation setting.} \\
\noindent
\textbf{Reply 3:} {We thank the Reviewer for this suggestion. Every reported result now comes from multiple independent runs, five seeds for the main configurations and three for the sweeps, with the mean and the standard deviation for every entry and the number of seeds stated in every table.

We gave particular attention to the two cases the Reviewer highlights. For the convergence-speed comparison, the single-seed relative metric has been replaced with rounds to fixed absolute accuracy, reported as a mean and a standard deviation over seeds and supported by a paired Wilcoxon signed-rank test over matched seeds. For the sparse participation setting, the case $K = 1$ at $N = 30$ now uses multiple seeds and reads $60.61 \pm 1.68$ for SCOPE-FD, and the participation sweep across $K \in \{1, 3, 5, 10\}$ is reported in the same way. Please refer to Section~VI-A for the protocol and to Section~VI-E for the sparse-participation sweep of the revised manuscript.}\\

\noindent
\textbf{Comment 4 (Ablation study):} \textit{The proposed scoring function consists of three complementary components. An ablation study comparing participation debt only, debt plus under-prediction bonus, debt plus class-coverage penalty, and the complete SCOPE-FD score would help quantify the contribution of each component and provide additional insight into the design choices.} \\
\noindent
\textbf{Reply 4:} {We thank the Reviewer for this suggestion, which matches the first comment of Reviewer 2. We have added the exact four-way comparison. Over five seeds at the headline configuration, participation debt only reaches $71.26 \pm 1.14$, debt plus the under-prediction bonus $70.55 \pm 1.19$, debt plus the class-coverage penalty $71.56 \pm 0.98$ and the complete score $71.21 \pm 0.99$, with a participation Gini of $1.33\%$ for all four.

The ablation shows that the participation-debt term alone provides the entire fairness guarantee, and that the contribution of the two auxiliary terms is confined to convergence speed, where the complete score reaches sixty percent absolute accuracy in $13.0 \pm 2.7$ rounds against $14.0 \pm 2.7$ for debt only. Accuracy is evaluated every five rounds, so that column is quantized and its dispersion is now printed rather than left implicit. The claims in the manuscript have been adjusted to match. Please refer to Section~VI-B and Table~II of the revised manuscript.}\\

\noindent
\textbf{Comment 5 (Privacy discussion):} \textit{The proposed method requires each client to upload its normalized label histogram. Although the manuscript briefly discusses possible privacy-preserving alternatives, a slightly more detailed discussion of the practical implications of sharing these statistics would further improve the completeness of the paper.} \\
\noindent
\textbf{Reply 5:} {We thank the Reviewer for this comment, which matches the second comment of Reviewer 2. We have replaced the brief discussion with a measured study and an explicit statement of what is disclosed.

What the server learns is the normalized label histogram of each client, a $C$-dimensional vector of class proportions. It does not learn any sample, any label of any sample, or the size of the local dataset, and the disclosure is static, since the histogram is uploaded once rather than per round.

We then measured what it costs to remove that disclosure. Against $71.99 \pm 0.63$ for the undefended selector, a Laplace mechanism on the histogram reaches $72.01 \pm 0.86$ at $\varepsilon_{\mathrm{dp}} = 0.1$, $71.71 \pm 0.72$ at $\varepsilon = 0.5$, $71.97 \pm 0.69$ at $\varepsilon = 1$, $71.78 \pm 0.80$ at $\varepsilon = 2$ and $71.85 \pm 0.73$ at $\varepsilon_{\mathrm{dp}} = 5$, and a server-side surrogate formed from the public-set logits reaches $71.93 \pm 0.67$. Every difference lies inside one standard deviation and the participation Gini is $1.33\%$ for all seven variants.

The practical implication is that the histogram is optional rather than central. It feeds only the under-prediction bonus and the class-coverage penalty, which the four-way ablation shows do not change final accuracy, while the participation-debt term that carries the fairness guarantee never reads it. A deployment that regards the disclosure as unacceptable can adopt the surrogate, or the debt-only variant, at no measured cost. Please refer to Section~VI-H of the revised manuscript.}\\

\end{document}
